% context: 9th-grade physics classwork, Unit 2 kinematics, Giancoli 2-7 (pp. 33-37)
% topic: free fall with a = g = 9.80 m/s^2, the four constant-acceleration equations
%        with y in place of x; dropped from rest (down positive), distance grows as
%        t^2, Galileo's heavy-and-light claim, a ball thrown upward (up positive):
%        v = 0 but a = -g at the top, symmetry of the flight, air resistance
% structure: two pages; page 1 the equations, a dropped ball table (Example 2-10)
%        and the Leaning Tower of Pisa; page 2 a ball thrown upward (Examples 2-12,
%        2-13, 2-14, 2-15) and a raindrop with and without air.
%        No quadratic formula (Example 2-16)
% data: Leaning Tower of Pisa 55.86 m on its low side; raindrop terminal speed about
%        9 m/s (contexts.md #8). The cloud height and throw speed are chosen values
\documentclass[12pt, twoside]{article}
%\documentclass[12pt, twoside]{article}
\usepackage[letterpaper, margin=1in, headsep=0.2in]{geometry}
\setlength{\headheight}{0.6in}
%\usepackage[english]{babel}
\usepackage[utf8]{inputenc}
\usepackage{microtype}
\usepackage{amsmath}
\usepackage{amssymb}
%\usepackage{amsfonts}
\usepackage[nomessages]{fp} %\FPeval{\var-name}{2*sin(pi/6)}
\usepackage{siunitx} %units in math. eg 20\milli\meter
\usepackage{yhmath} % for arcs, overparenth command
\usepackage{tikz} %graphics
\usetikzlibrary{quotes, angles, arrows, arrows.meta}
\usepackage{graphicx} %consider setting \graphicspath{{images/}}
\usepackage{parskip} %no paragraph indent
\usepackage{enumitem}
\usepackage{multicol}
\usepackage{venndiagram}

\usepackage{fancyhdr}
\pagestyle{fancy}
\fancyhf{}
\renewcommand{\headrulewidth}{0pt} % disable the underline of the header
\raggedbottom
\hfuzz=2mm %suppresses overfull box warnings

\usepackage{hyperref}
\usepackage{float}

\title{Physics}
\author{Chris Huson}
\date{October 2026}

\fancyhead[LE]{\thepage}
\fancyhead[RO]{\thepage \\ Nickname: \hspace{2.25cm} \,\\ \,}
\fancyhead[LO]{La Scuola d'Italia / Huson / Physics \\* 30 October 2026}

\begin{document}

\subsubsection*{2.5 Classwork: Falling Objects}

\textbf{Reference (provided).} \quad For a falling object, the acceleration is
$g = 9.80$ m/s$^2$, downward.

\vspace{0.1cm}
{\hspace{0.5cm}%
\renewcommand{\arraystretch}{1.8}%
\begin{tabular}{l @{\hspace{2cm}} l}
  $v = v_0 + at$ & $v^2 = v_0^2 + 2a(y - y_0)$ \\
  $y = y_0 + v_0 t + \frac{1}{2}at^2$ & $\bar{v} = \dfrac{v + v_0}{2}$
\end{tabular}}

\vspace{0.2cm}
Ignore air resistance unless the problem says otherwise. ``Dropped'' means
$v_0 = 0$. Choose which direction is positive, up or down, and keep it for the whole
problem. List what you know and what you want, then choose the equation.

\vspace{0.2cm}

\begin{enumerate}[itemsep=0.4cm]

\item A ball is dropped from a tall building. Take \emph{down} as positive, so
  $a = +9.80$ m/s$^2$, and $y_0 = 0$.
  \begin{enumerate}[itemsep=0.4cm]
    \item Fill in the table.

      \vspace{0.1cm}
      \renewcommand{\arraystretch}{1.5}
      \begin{tabular}{| c | p{1.6cm} | p{1.6cm} | p{1.6cm} | p{1.6cm} |}
        \hline
        $t$ (s) & \centering 0 & \centering 1.00 & \centering 2.00 &
          \centering\arraybackslash 3.00 \\ \hline
        $v$ (m/s) & & & & \\ \hline
        $y$ (m) & & & & \\ \hline
      \end{tabular}
    \item At 2.00 s, the ball has fallen how many times as far as at 1.00 s? At
      3.00 s? What pattern do you see?
      \vspace{1.2cm}
  \end{enumerate}

\item Galileo is said to have dropped balls from the Leaning Tower of Pisa, which is
  55.86 m tall on its low side. Take down as positive.
  \begin{enumerate}[itemsep=0.4cm]
    \item How long does a ball dropped from the top take to reach the ground?
      (Solve for $t^2$ first, then take the square root.)
      \vspace{2.2cm}
    \item How fast is it going when it lands?
      \vspace{1.6cm}
    \item A heavy iron ball and a light wooden ball of the same size are dropped
      together. Which lands first?
  \end{enumerate}

\newpage

\item A student throws a ball straight up. It leaves the hand at 14.7 m/s. Take
  \emph{up} as positive, so $v_0 = +14.7$ m/s and $a = -9.80$ m/s$^2$.
  \begin{enumerate}[itemsep=0.4cm]
    \item At the highest point, what is the ball's velocity? What is its
      acceleration?
      \vspace{1.4cm}
    \item How long does the ball take to reach the highest point?
      \vspace{2.4cm}
    \item How high above the hand does it go?
      \vspace{2.4cm}
    \item How long after the throw does the ball come back to the hand, and what
      is its velocity then?
      \vspace{2.4cm}
  \end{enumerate}

\item A raindrop falls from a cloud 1000 m above the ground. Take down as positive.
  \begin{enumerate}[itemsep=0.4cm]
    \item If there were no air, how fast would the raindrop be going when it reached
      the ground?
      \vspace{2.6cm}
    \item Real raindrops land at about 9 m/s. Why is that so much slower?
  \end{enumerate}

\end{enumerate}
\end{document}
